\subsection{The open-orbit method has a precise higher-height limit}
\begin{finding}[\fid{F29}: the same flag-orbit mechanism stops at height two]\label{f:sphericity}
Let $n\ge3$, $E=\C^n\otimes\C^n$, $N=n^2$, and
$G=SL_n\times SL_n$. For every $2\le r\le N$, $G$ has no
open orbit on $\operatorname{Flag}(1,\ldots,r;E)$.
Equivalently, $G\times B_r$ has no open orbit on $E\otimes\C^r$,
where $B_r$ is a Borel subgroup. For $r=2$ the generic orbit has
codimension $n-2$. The same open-orbit proof underlying Manivel's
height-two formula therefore does not extend to these heights,
even after restricting the output length to two.
\end{finding}
\begin{proof}
On the dense locus of injections $(\C^r)^*\to E$, quotienting
by the appropriate dual Borel records the flag of successive images.
For each flag the compatible ordered injections form one free Borel
orbit. This is the usual frame bundle of the filtered tautological
bundle, locally a principal Borel bundle. It proves the equivalence
of the two open-orbit questions. The partial flag dimension is
\[
 rN-r(r+1)/2,\qquad \dim G=2(N-1).
\]
At $r=3$ the difference is $N-4>0$; successive increments of the
flag dimension are $N-r-1\ge0$ for $3\le r<N$. Thus no open
orbit exists for $r\ge3$. At full length the deficit is
$(N-1)(N-4)/2$.

For $r=2$, view $E$ as matrices, with action $X\mapsto gXh^t$.
The invertible first lines form a dense open orbit in $\mathbf P(E)$:
over $\C$ a scalar can normalize any invertible matrix's determinant.
A generic flag can therefore be written
$\C I\subset\C I+\C B$, with $B$ regular semisimple.
Its stabilizer satisfies $gh^t=cI$, $c^n=1$, and
$gBg^{-1}=aI+bB$, $b\ne0$. The affine transformation
$x\mapsto a+bx$ must permute the distinct eigenvalues of $B$.
Only finitely many such affine transformations exist, because the
images of two distinct eigenvalues determine $a,b$. The identity
component therefore has $c=1,a=0,b=1$ and is the centralizer of
$B$ in $SL_n$, a torus of dimension $n-1$. The generic orbit has
dimension $2N-n-1$ in a flag variety of dimension $2N-3$,
giving codimension $n-2>0$. Any open orbit would meet this dense
regular-semisimple locus, which is impossible.
\end{proof}
For $r=1$ the invertible matrix lines do form an open orbit, so the
boundary is real. The precise source premise being tested is Manivel
\cite[Lemma~1 and Corollary~1, \S3, p.~3]{Manivel}. Failure
of this open-orbit premise excludes its multiplicity-free algebra
mechanism; it does not exclude positive bases, useful generators,
or counting rules for multiplicity-bearing invariant algebras.
